% Conclusion restricted to the current route-only findings registry.

\section{Conclusion}
\label{sec:conclusion}

Realised trading paths make vehicle-currency use directly observable. They separate the assets that ever intermediate a trade from the smaller set that carries most routed activity. In our sample, stablecoins become the leading vehicles by value relative to native assets, although the transition reverses more than once. USDT accounts for most of the increase within the stablecoin category.

Where trading occurs is central to this transition. Most of the increase comes from endpoint pairs entering or leaving, changes in activity among continuing pairs, and greater use of indirect routes. Direct stable-for-native switching inside continuing markets contributes little. A model that holds the set of markets fixed therefore captures only part of the formation of vehicle dominance.

The usual liquidity feedback begins with currency choice inside an established market: greater use deepens liquidity and makes future use more attractive. The route evidence adds a market-formation margin. Changes in which pairs trade and which of them route indirectly can extend the pools and software support associated with a currency even when intermediary choice inside a continuing market changes little.

Liquidity provision, venue reach, endpoint demand, and routing software may explain why activity moved toward USDT-centred routes; realised paths alone cannot separate these mechanisms. What the paths establish is that vehicle dominance is shaped by the formation of routed markets as well as by the choice of intermediary within markets already in place.
